\chapter{Stochastic Processes, SDEs, and the Ornstein--Uhlenbeck Channel}
\label{ch:stochproc}\label{ch:sde}

Chapter~\ref{ch:statmech} answered the question \emph{``which
distribution?''} --- Boltzmann's. This chapter answers \emph{``how does
probability move?''}, in discrete and in continuous time, and it earns its
place in the thesis three times over. Markov chains are the \emph{data}
whose diffusion score we compute exactly (the AR(1) sequence of
Section~\ref{sec:ar1} is the running example of every research chapter);
stochastic differential equations are the \emph{corruption and generation
processes} of diffusion models (the Ornstein--Uhlenbeck channel of
Section~\ref{sec:ou} appears on every page of
Chapters~\ref{ch:gaussian}--\ref{ch:bp}); and the Fokker--Planck and
time-reversal theorems of Sections~\ref{sec:fokker-planck}--\ref{sec:reversal}
are the mathematical licence for the entire generative programme reviewed in
Chapter~\ref{ch:diffusion}. Everything proved here is used later; nothing is
included for completeness' sake. Standard references are
\citet{robert2004monte, brockwell1991time, oksendal2003stochastic,
gardiner2009stochastic, risken1996fokker}.

\section{Markov chains: kernels, stationarity, reversibility}
\label{sec:markov}

\begin{definition}[Markov chain]
A sequence of random variables $(X_k)_{k \geq 0}$ with values in a state
space $\mathcal{X}$ is a \emph{Markov chain} if for every $k$ and every
measurable $A \subseteq \mathcal{X}$,
\begin{equation}
  \Prob\big(X_{k+1} \in A \,\big|\, X_k, X_{k-1}, \dots, X_0\big)
  \;=\; \Prob\big(X_{k+1} \in A \,\big|\, X_k\big) .
  \label{eq:markov-property}
\end{equation}
The conditional law is encoded by a \emph{transition kernel}
$M(x' \,|\, x)$: a density (or matrix, for discrete $\mathcal{X}$) with
$\int M(x'|x)\,\dd x' = 1$ for every $x$.
\end{definition}

The Markov property \eqref{eq:markov-property} is exactly a statement about
\emph{factorisation}: the joint law of the first $K$ states is
\begin{equation}
  p(x_0, x_1, \dots, x_{K-1})
  \;=\; p_0(x_0)\, \prod_{k=0}^{K-2} M(x_{k+1} \,|\, x_k) .
  \label{eq:chain-factorisation}
\end{equation}
This one-line factorisation is the single most important equation of the
thesis. It says that the joint distribution of a Markov trajectory is a
product of \emph{local} terms, each touching two consecutive variables ---
the structure that makes the chain a \emph{tree} in the factor-graph sense
of Chapter~\ref{ch:bp}, that makes its precision matrix tridiagonal in the
Gaussian case of Chapter~\ref{ch:gaussian}, and that makes belief
propagation exact.

Marginal distributions evolve by integrating the kernel:
$p_{k+1}(x') = \int M(x'|x)\, p_k(x)\,\dd x$, compactly $p_{k+1} = M p_k$.
A distribution $\pi$ is \emph{stationary} (invariant) for $M$ if
$\pi = M\pi$; a chain started in $\pi$ stays in $\pi$ forever, and the
trajectory is then a \emph{stationary process}. The kernel satisfies
\emph{detailed balance} with respect to $\pi$ if
\begin{equation}
  \pi(x)\, M(x' \,|\, x) \;=\; \pi(x')\, M(x \,|\, x')
  \qquad \text{for all } x, x' :
  \label{eq:detailed-balance}
\end{equation}
the equilibrium probability flux from $x$ to $x'$ equals the reverse flux,
so the chain, watched in equilibrium, is statistically indistinguishable run
forwards or backwards. Detailed balance implies stationarity (integrate
\eqref{eq:detailed-balance} over $x$), and its power is practical: it is a
\emph{local, pointwise} condition, checkable without computing any integral.

Stationarity alone does not guarantee that the chain \emph{finds} $\pi$; for
that one needs the chain to be able to reach everywhere
(\emph{irreducibility}) without getting trapped in cycles
(\emph{aperiodicity}). Under these conditions the fundamental convergence
theorem holds \citep{robert2004monte}: the marginal law of $X_k$ converges
to $\pi$ from any initial condition, and time averages converge to
expectations under $\pi$ (the ergodic theorem),
\begin{equation}
  \frac{1}{N} \sum_{k=1}^{N} f(X_k)
  \;\xrightarrow[N \to \infty]{\text{a.s.}}\;
  \E_\pi[f] .
  \label{eq:ergodic}
\end{equation}
Equation \eqref{eq:ergodic} is the rigorous form of the ergodic postulate
invoked in Section~\ref{sec:ensembles}, with the physical dynamics replaced
by an algorithmic one.

\section{The AR(1) chain: the running example of the thesis}\label{sec:ar1}

We now introduce the specific Markov chain that the research chapters solve
completely: the first-order autoregressive (AR(1)) process, the canonical
Gaussian Markov chain \citep{brockwell1991time}. Let
$\alpha \in (-1, 1)$ and let $(\eta_k)$ be i.i.d.\ $\Nd(0, \sigma_\eta^2)$
innovations; define
\begin{equation}
  a_{k+1} \;=\; \alpha\, a_k + \eta_k ,
  \qquad k = 0, 1, \dots, K-2 .
  \label{eq:ar1}
\end{equation}
The transition kernel is Gaussian,
$M(a' | a) = \Nd(a'; \alpha a, \sigma_\eta^2)$, and the joint law of the
trajectory follows from \eqref{eq:chain-factorisation}. Everything about
this chain is computable in closed form; the next toolbox derives the three
facts used throughout.

\begin{toolbox}[tb:ar1]{Stationary distribution and autocovariance of AR(1)}
\emph{(i) Stationary variance.} Suppose $a_k \sim \Nd(0, \sigma^2)$ and
impose that $a_{k+1}$ has the same variance. Since $\eta_k$ is independent
of $a_k$,
\begin{equation*}
  \operatorname{Var}(a_{k+1})
  = \alpha^2 \operatorname{Var}(a_k) + \sigma_\eta^2
  \;\stackrel{!}{=}\; \operatorname{Var}(a_k)
  \quad\Longrightarrow\quad
  \sigma_\infty^2 = \frac{\sigma_\eta^2}{1 - \alpha^2}.
\end{equation*}
The fixed point exists precisely because $|\alpha| < 1$; the map
$\sigma^2 \mapsto \alpha^2 \sigma^2 + \sigma_\eta^2$ is a contraction, so
\emph{any} initial variance converges to $\sigma_\infty^2$ geometrically.
Since Gaussianity is preserved by the affine recursion, the stationary law
is $\pi = \Nd(0, \sigma_\infty^2)$.

\emph{(ii) Autocovariance.} In the stationary regime, for $d \geq 1$, unroll
the recursion $d$ times:
$a_{k+d} = \alpha^d a_k + \sum_{j=0}^{d-1} \alpha^{j} \eta_{k+d-1-j}$.
All innovation terms are independent of $a_k$, hence
\begin{equation*}
  \operatorname{Cov}(a_k, a_{k+d})
  = \alpha^d\, \operatorname{Var}(a_k)
  = \alpha^d\, \sigma_\infty^2 ,
\end{equation*}
and by symmetry $\operatorname{Cov}(a_i, a_j) = \alpha^{|i-j|}
\sigma_\infty^2$ for all $i, j$: correlations decay \emph{geometrically}
with the lag, with rate $\alpha$. The covariance matrix of $K$ consecutive
frames,
\begin{equation*}
  (\Sigz)_{ij} = \sigma_\infty^2\, \alpha^{|i-j|},
\end{equation*}
depends on $(i,j)$ only through $i - j$: it is a symmetric \emph{Toeplitz}
matrix, a structure we exploit heavily in Chapter~\ref{ch:gaussian}.

\emph{(iii) Detailed balance.} With $\pi = \Nd(0, \sigma_\infty^2)$ and
$M(a'|a) = \Nd(a'; \alpha a, \sigma_\eta^2)$, both sides of
\eqref{eq:detailed-balance} are proportional to
\begin{equation*}
  \exp\!\left[
    -\frac{a^2}{2\sigma_\infty^2}
    -\frac{(a' - \alpha a)^2}{2\sigma_\eta^2}
  \right]
  = \exp\!\left[
    -\frac{a^2 + (a')^2 - 2\alpha\, a a'}{2\sigma_\eta^2}
  \right],
\end{equation*}
where the exponent on the right is \emph{symmetric} in $(a, a')$. Hence the
stationary AR(1) chain is reversible: statistically, the movie of the chain
looks the same played in either direction. \hfill$\square$
\end{toolbox}

The AR(1) chain is thus the perfect miniature of everything this chapter
discusses: a Markov kernel with an explicit stationary law, geometric
convergence, reversibility, and --- decisive for us --- \emph{exact Gaussian
algebra} at every step. When in Chapter~\ref{ch:gaussian} we corrupt each
frame of this chain with independent noise and ask for the joint score of
the result, every quantity will remain in closed form.

\section{Sampling by dynamics: MCMC and Langevin}\label{sec:mcmc}

Sampling inverts the logic of the previous section. There, a kernel $M$ was
given and we asked for its stationary law $\pi$. In sampling problems we are
given $\pi$ --- typically a Boltzmann distribution
$\pi(x) \propto e^{-\beta H(x)}$ known only up to its partition function ---
and we must \emph{design} a kernel whose stationary law is $\pi$; the
ergodic theorem \eqref{eq:ergodic} then converts one long trajectory into
expectation values. This is how the negative phase of Boltzmann-machine
learning \eqref{eq:bm-learning} is estimated in practice, and the idea
originates, fittingly, in computational statistical mechanics
\citep{metropolis1953equation}.

The \emph{Metropolis--Hastings} chain \citep{metropolis1953equation,
hastings1970monte} moves from $x$ by proposing $x' \sim g(\cdot|x)$ from any
tractable proposal kernel and accepting the move with probability
\begin{equation}
  A(x \to x') \;=\; \min\!\left\{ 1,\;
    \frac{\pi(x')\, g(x \,|\, x')}{\pi(x)\, g(x' \,|\, x)}
  \right\},
  \label{eq:mh-acceptance}
\end{equation}
staying at $x$ otherwise. The target enters only through the \emph{ratio}
$\pi(x')/\pi(x)$, so the intractable partition function cancels: MH samples
from $e^{-\beta H}/Z$ while only ever evaluating energy differences. The
acceptance rule is engineered so that the resulting kernel satisfies
detailed balance \eqref{eq:detailed-balance} with respect to $\pi$; the
short verification is Derivation~\ref{drv:mh} in
Appendix~\ref{app:derivations}. When $x$ is high-dimensional but each
conditional $\pi(x_i | x_{\setminus i})$ is tractable, the \emph{Gibbs
sampler} \citep{geman1984stochastic} cycles through the coordinates
resampling each from its conditional --- an MH move whose acceptance
probability is identically one. Gibbs sampling thrives exactly when the
model has sparse conditional structure: for a Markov chain,
$\pi(x_k | x_{\setminus k})$ depends only on the two neighbours, and for the
bipartite RBM of Section~\ref{sec:boltzmann-machines} each layer updates in
one parallel step. The reader should hear, already, the theme of
Chapter~\ref{ch:bp}: \emph{inference cost follows graph structure}.

For continuous $x$ and differentiable energy, a different design principle
uses gradients. Consider the discrete-time update
\begin{equation}
  x_{k+1} \;=\; x_k + \epsilon\, \nabla \log \pi(x_k)
  \;+\; \sqrt{2\epsilon}\; \xi_k ,
  \qquad \xi_k \sim \Nd(0, I),
  \label{eq:ula}
\end{equation}
the (unadjusted) \emph{Langevin algorithm}. The drift pushes towards high
probability; the noise maintains exploration. Two observations preview what
follows. First, \eqref{eq:ula} requires only the \emph{score}
$\nabla \log \pi$ --- the normalising constant is never needed; this is
precisely why the score is the natural object for generative modelling
(Chapter~\ref{ch:diffusion}), and why so much of this thesis is devoted to
computing scores exactly. Second, as $\epsilon \to 0$ the recursion
\eqref{eq:ula} becomes a stochastic differential equation --- the Langevin
diffusion --- whose stationary law is exactly $\pi$. Making that limit
precise requires the It\^o calculus and Fokker--Planck theory that occupy
the rest of the chapter, where the claim is verified by direct computation.

\section{From ODEs to SDEs: the It\^o calculus}\label{sec:ode-to-sde}

An ordinary differential equation $\dot x = f(x, t)$ moves a point along a
deterministic velocity field. Physical systems, however, feel rapid,
unresolved forces --- collisions, thermal agitation --- whose cumulative
effect is well modelled by adding noise:
\begin{equation}
  \dd X_t \;=\; \underbrace{f(X_t, t)\,\dd t}_{\text{drift}}
  \;+\; \underbrace{g(t)\,\dd W_t}_{\text{diffusion}} .
  \label{eq:generic-sde}
\end{equation}
Here $W_t$ is a \emph{Wiener process} (standard Brownian motion): the unique
process with $W_0 = 0$, independent increments,
$W_t - W_s \sim \Nd(0, t-s)$, and continuous paths. Equation
\eqref{eq:generic-sde} is shorthand for the integral equation
$X_t = X_0 + \int_0^t f(X_s, s)\,\dd s + \int_0^t g(s)\,\dd W_s$, and the
whole subtlety lives in the last term: Brownian paths are nowhere
differentiable, so $\int g\,\dd W$ cannot be a Riemann--Stieltjes integral.
It\^o's construction \citep{ito1944stochastic} defines it as the limit of
left-endpoint sums,
\begin{equation}
  \int_0^t g(s)\, \dd W_s
  \;=\; \lim_{n \to \infty} \sum_{j=0}^{n-1}
        g(s_j)\,\big( W_{s_{j+1}} - W_{s_j} \big),
  \label{eq:ito-integral}
\end{equation}
where the evaluation at the \emph{left} endpoint $s_j$ makes the integrand
independent of the coming increment. Two consequences follow immediately:
the It\^o integral is a martingale (zero mean), and its variance is given by
the \emph{It\^o isometry},
$\E\big[ (\int_0^t g\,\dd W)^2 \big] = \int_0^t g(s)^2\,\dd s$, both because
cross terms $\E[g(s_j)(W_{s_{j+1}}-W_{s_j})] $ vanish by independence.

The price of the construction is a new chain rule. Its origin is the
\emph{quadratic variation} of Brownian motion: increments scale as
$\dd W_t \sim \sqrt{\dd t}$, so second-order terms in $\dd W$ are
first-order in $\dd t$ and cannot be discarded.

\begin{toolbox}{It\^o's lemma, and where the extra term comes from}
Let $Y_t = \phi(X_t, t)$ with $\phi$ twice differentiable and $X_t$ solving
\eqref{eq:generic-sde}. Taylor-expand $\phi$ to second order:
\begin{equation*}
  \dd Y
  = \frac{\partial \phi}{\partial t}\,\dd t
  + \frac{\partial \phi}{\partial x}\,\dd X
  + \frac{1}{2} \frac{\partial^2 \phi}{\partial x^2}\,(\dd X)^2
  + \dots
\end{equation*}
Substitute $\dd X = f\,\dd t + g\,\dd W$ and expand the square:
\begin{equation*}
  (\dd X)^2 = f^2 (\dd t)^2 + 2 f g\, \dd t\, \dd W + g^2 (\dd W)^2 .
\end{equation*}
Now use the multiplication table of It\^o calculus, which is nothing but the
scaling $\dd W \sim \sqrt{\dd t}$ made systematic:
\begin{equation*}
  (\dd t)^2 = 0, \qquad \dd t\, \dd W = 0, \qquad (\dd W)^2 = \dd t .
\end{equation*}
The last identity is the rigorous statement that the quadratic variation of
$W$ over $[0,t]$ equals $t$: for a partition of mesh $\to 0$,
$\sum_j (W_{s_{j+1}} - W_{s_j})^2 \to t$ in mean square, since the sum has
mean $\sum_j (s_{j+1} - s_j) = t$ and variance
$2\sum_j (s_{j+1}-s_j)^2 \to 0$. Keeping the surviving terms:
\begin{equation*}
  \dd Y = \left(
    \frac{\partial \phi}{\partial t}
    + f\,\frac{\partial \phi}{\partial x}
    + \frac{g^2}{2}\,\frac{\partial^2 \phi}{\partial x^2}
  \right) \dd t
  \;+\; g\,\frac{\partial \phi}{\partial x}\; \dd W .
\end{equation*}
The half-Laplacian term --- absent from the classical chain rule --- is the
engine of everything that follows: it is why densities \emph{spread}, why
the Fokker--Planck equation has a second-order term, and why
$\exp$-transformations of SDEs pick up variance corrections.
\hfill$\square$
\end{toolbox}

\section{The Ornstein--Uhlenbeck process: the corruption channel}
\label{sec:ou}

The noising process used by diffusion models --- and by every result of this
thesis --- is the \emph{Ornstein--Uhlenbeck process}
\citep{uhlenbeck1930theory}: Brownian motion pulled back to the origin by a
linear spring,
\begin{equation}
  \dd X_t \;=\; -X_t\,\dd t \;+\; \sqrt{2}\;\dd W_t ,
  \qquad X_0 = a .
  \label{eq:ou-sde}
\end{equation}
(The drift and diffusion constants are a normalisation choice --- this is
the ``variance-preserving'' convention of the diffusion literature
\citep{song2021scorebased}; any other choice rescales time.) The OU process
is the unique Gaussian, Markov, stationary process in continuous time, and
it is exactly solvable.

\begin{toolbox}{Exact solution of the OU SDE by integrating factor}
Multiply \eqref{eq:ou-sde} by $e^{t}$ and consider
$Y_t = e^{t} X_t$. By It\^o's lemma with $\phi(x,t) = e^{t}x$ (here
$\partial^2 \phi/\partial x^2 = 0$, so no correction term arises):
\begin{equation*}
  \dd Y_t
  = e^{t} X_t\,\dd t + e^{t}\big( -X_t\,\dd t + \sqrt{2}\,\dd W_t \big)
  = \sqrt{2}\, e^{t}\, \dd W_t .
\end{equation*}
The drift cancels exactly --- that is what the integrating factor is for.
Integrate from $0$ to $t$ and undo the substitution:
\begin{equation*}
  X_t \;=\; a\, e^{-t} \;+\; \sqrt{2} \int_0^t e^{-(t-s)}\, \dd W_s .
\end{equation*}
The integral is Gaussian (limit of Gaussian sums), with mean zero and
variance given by the It\^o isometry:
\begin{equation*}
  \operatorname{Var}(X_t)
  = 2 \int_0^t e^{-2(t-s)}\, \dd s
  = 2\, e^{-2t}\, \frac{e^{2t} - 1}{2}
  = 1 - e^{-2t} .
\end{equation*}
Hence the \emph{transition kernel} of the OU channel is
\begin{equation*}
  p(x_t \,|\, X_0 = a)
  \;=\; \Nd\!\big( x_t;\; a\,e^{-t},\; \Deltat \big),
  \qquad
  \boxed{\;\Deltat \;=\; 1 - e^{-2t}\;} .
\end{equation*}
\hfill$\square$
\end{toolbox}

The boxed kernel is used on every page of Chapters~\ref{ch:gaussian}
--\ref{ch:bp}, so we fix its reading now. Corrupting a data point $a$ for a
time $t$ produces
\begin{equation}
  x \;=\; a\, e^{-t} \;+\; \sqrt{\Deltat}\;\varepsilon,
  \qquad \varepsilon \sim \Nd(0, 1):
  \label{eq:ou-channel}
\end{equation}
the signal is \emph{contracted} by $e^{-t}$ while isotropic noise of
variance $\Deltat$ fills in. At $t = 0$: $x = a$ (no corruption); as
$t \to \infty$: $x \sim \Nd(0, 1)$ regardless of $a$ (all information
destroyed, the equilibrium of the spring). The signal-to-noise ratio
$e^{-2t}/\Deltat$ sweeps continuously from $\infty$ to $0$: the OU channel
interpolates between the data distribution and a standard Gaussian, which is
precisely what a diffusion model needs.

Two structural remarks, both load-bearing later. First, the OU kernel is
Gaussian \emph{in $a$ as well as in $x$}: viewed as a function of the clean
variable, $\Nd(x; a e^{-t}, \Deltat) \propto
\exp[-(e^{-2t}/2\Deltat)(a - x e^{t})^2]$ --- a Gaussian ``likelihood''
factor. This is the fact that makes the posteriors of
Chapter~\ref{ch:gaussian} jointly Gaussian and the messages of
Chapter~\ref{ch:bp} closable. Second, if a random vector
$a \sim \Nd(\mu, \Sigma_0)$ is pushed through \eqref{eq:ou-channel}
coordinate-wise with \emph{independent} noises, the output is Gaussian with
\begin{equation}
  \mu_t = e^{-t}\mu,
  \qquad
  \Sigt \;=\; e^{-2t}\, \Sigma_0 \;+\; \Deltat\, I ,
  \label{eq:ou-cov-map}
\end{equation}
since means and covariances transform linearly and the added noise is white.
Equation \eqref{eq:ou-cov-map} --- shrink the covariance, add a multiple of
the identity --- is the entire ``forward process'' of this thesis in one
line.

\section{The Fokker--Planck equation: transporting densities}
\label{sec:fokker-planck}

The SDE \eqref{eq:generic-sde} moves \emph{samples}; the corresponding
motion of the \emph{density} $p_t(x)$ is governed by the Fokker--Planck (or
forward Kolmogorov) equation. This change of viewpoint --- from trajectories
to densities --- is the same one Liouville's theorem performed for
Hamiltonian flow, and the derivation is parallel, with the It\^o correction
supplying a new, diffusive term.

\begin{toolbox}{Derivation of the Fokker--Planck equation}
Let $\phi$ be an arbitrary smooth test function of compact support, and
compute $\frac{\dd}{\dd t}\E[\phi(X_t)]$ in two ways.

\emph{Way 1: through the samples.} By It\^o's lemma,
\begin{equation*}
  \dd \phi(X_t)
  = \Big( f\,\phi' + \tfrac{g^2}{2}\,\phi'' \Big)\dd t
  + g\,\phi'\,\dd W_t ,
\end{equation*}
and taking expectations kills the martingale term:
\begin{equation*}
  \frac{\dd}{\dd t}\,\E[\phi(X_t)]
  = \E\Big[ f(X_t,t)\,\phi'(X_t) \Big]
  + \frac{g(t)^2}{2}\, \E\big[ \phi''(X_t) \big] .
\end{equation*}

\emph{Way 2: through the density.} $\E[\phi(X_t)] = \int \phi(x)\,
p_t(x)\,\dd x$, so
\begin{equation*}
  \frac{\dd}{\dd t}\,\E[\phi(X_t)]
  = \int \phi(x)\, \frac{\partial p_t(x)}{\partial t}\, \dd x .
\end{equation*}

Equate the two, and integrate Way 1 by parts to move all derivatives off
$\phi$ (boundary terms vanish by compact support): once for the drift term,
twice for the diffusion term,
\begin{equation*}
  \int \phi\, \frac{\partial p_t}{\partial t}
  = \int \phi \left[
      -\frac{\partial}{\partial x}\big( f\, p_t \big)
      + \frac{g^2}{2}\, \frac{\partial^2 p_t}{\partial x^2}
    \right] \dd x .
\end{equation*}
Since $\phi$ is arbitrary, the integrands must agree pointwise:
\begin{equation*}
  \boxed{\;
  \frac{\partial p_t}{\partial t}
  = -\frac{\partial}{\partial x}\big( f(x,t)\, p_t(x) \big)
  + \frac{g(t)^2}{2}\, \frac{\partial^2 p_t(x)}{\partial x^2} .
  \;}
\end{equation*}
The first term transports probability along the drift (as in Liouville);
the second spreads it diffusively --- it is the macroscopic footprint of the
It\^o correction. \hfill$\square$
\end{toolbox}

\subsection{Stationary law of Langevin dynamics: closing the loop}

Apply the Fokker--Planck equation to the \emph{Langevin SDE} associated with
a potential $U$ (the continuous-time limit of the algorithm
\eqref{eq:ula}):
\begin{equation}
  \dd X_t = -\nabla U(X_t)\,\dd t + \sqrt{2\beta^{-1}}\;\dd W_t .
  \label{eq:langevin-sde}
\end{equation}
A stationary density solves
$0 = \nabla \cdot ( \nabla U\, p + \beta^{-1} \nabla p )$, i.e.\ the
probability current $J = -(\nabla U)\,p - \beta^{-1}\nabla p$ is
divergence-free; imposing the natural (decaying) boundary conditions forces
$J \equiv 0$, hence $\nabla p / p = -\beta \nabla U$, hence
\begin{equation}
  p_\infty(x) \;\propto\; e^{-\beta U(x)} :
  \label{eq:langevin-boltzmann}
\end{equation}
\emph{Langevin dynamics equilibrates to the Boltzmann distribution} of
Chapter~\ref{ch:statmech}. This is the two-line proof of the claim left
pending in Section~\ref{sec:mcmc}, and it explains the deep kinship of
sampling and physics: a Langevin sampler \emph{is} an overdamped physical
system at temperature $1/\beta$. Note also that the drift of
\eqref{eq:langevin-sde} is (up to $\beta$) the \emph{score} of the target,
$-\nabla U = \nabla \log p_\infty$: dynamics that sample a distribution are
driven by its score.

For the OU process, $U(x) = x^2/2$ and $\beta = 1$: the stationary law is
$\Nd(0,1)$, confirming the $t \to \infty$ limit of Section~\ref{sec:ou}.

\section{Time reversal: the mother theorem of diffusion models}
\label{sec:reversal}

One last piece of SDE theory converts everything above into a generative
method. Suppose data $X_0 \sim p_0$ are corrupted by a forward SDE
\eqref{eq:generic-sde} up to time $T$, producing marginals $p_t$. Can the
corruption be undone --- can we find an SDE which, run from
$X_T \sim p_T$ backwards, reproduces the same marginals? The answer
\citep{anderson1982reverse} is yes, and the reverse drift needs exactly one
unknown object:
\begin{equation}
  \dd X_t \;=\;
  \Big[ f(X_t, t) \;-\; g(t)^2\, \nabla_x \log p_t(X_t) \Big]\,\dd t
  \;+\; g(t)\, \dd \bar W_t ,
  \label{eq:reverse-sde}
\end{equation}
where time runs backwards from $T$ to $0$ and $\bar W$ is a reverse-time
Wiener process. The correction to the drift is the \emph{score}
$\nabla_x \log p_t$ of the noised marginal --- the object defined in
\eqref{eq:intro-score} and computed exactly, for Markov-chain data,
in Chapters~\ref{ch:gaussian}--\ref{ch:bp}. A heuristic that makes
\eqref{eq:reverse-sde} plausible follows from the Fokker--Planck equation:
rewriting its diffusion term as a transport term,
\begin{equation*}
  \frac{g^2}{2} \frac{\partial^2 p_t}{\partial x^2}
  = \frac{\partial}{\partial x}\!\left(
      \frac{g^2}{2}\, p_t\, \frac{\partial \log p_t}{\partial x}
    \right),
\end{equation*}
one sees that the density evolution of \eqref{eq:generic-sde} is
\emph{identical} to that of the noise-free ODE with velocity
$f - \tfrac{g^2}{2} \nabla \log p_t$ (the \emph{probability-flow ODE} of
\citealp{song2021scorebased}); running that transport backwards, and
re-injecting noise consistently, yields \eqref{eq:reverse-sde}. The rigorous
statement and proof are in \citet{anderson1982reverse}.

\section{From algorithmic chains to data chains}

A closing change of perspective, to set up the research problem. In this
chapter Markov chains began as \emph{algorithms}: their stationary law was
the target, their trajectory a means to an end. From
Chapter~\ref{ch:gaussian} onwards the chain \eqref{eq:ar1} plays the
opposite role: it is the \emph{data-generating process}, its trajectory
$(a_0, \dots, a_{K-1})$ is one data point --- a ``video'' of $K$ frames ---
and the object of interest is the probability law of the \emph{whole
trajectory} after each frame has been independently corrupted by the OU
channel \eqref{eq:ou-channel}. The factorisation
\eqref{eq:chain-factorisation} then stops being a computational convenience
and becomes the \emph{structure to be discovered}: the entire research
programme of this thesis can be phrased as asking how much of
\eqref{eq:chain-factorisation} survives noising, and how an inference
algorithm --- or a neural network --- can exploit what survives. In one
sentence, the chapter's summary: the OU channel destroys data at a known
rate, the Fokker--Planck equation describes the destruction at the level of
densities, and Anderson's theorem says the destruction is invertible by
anyone who knows the score $\nabla \log p_t$ --- which is why the rest of
this thesis is about computing that score, exactly, when the data have
Markov structure.
